\documentclass[
10pt, % Main document font size
a4paper, % Paper type, use 'letterpaper' for US Letter paper
oneside, % One page layout (no page indentation)
%twoside, % Two page layout (page indentation for binding and different headers)
headinclude,footinclude, % Extra spacing for the header and footer
BCOR5mm, % Binding correction
]{scrartcl}
\input{stru/structure.tex}
\hyphenation{Fortran hy-phen-ation} 

\title{\normalfont{FTML practical session 6}} 
\author{\spacedlowsmallcaps{}
} 
\date{\today}
\begin{document}
\renewcommand{\sectionmark}[1]{\markright{\spacedlowsmallcaps{#1}}}
\lehead{\mbox{\llap{\small\thepage\kern1em\color{halfgray} \vline}\color{halfgray}\hspace{0.5em}\rightmark\hfil}} 
\pagestyle{scrheadings}
\maketitle 
\setcounter{tocdepth}{3} 

\section{\large\color{Blue}Line search for least squares}

We first note that 

\begin{equation}
    \begin{aligned}
        \label{eq:1}
	\nabla_{\theta}f(\alpha( \gamma)) &= H \alpha( \gamma) - \frac{1}{n} X^Ty\\
	&= H (\theta_{t}- \gamma \nabla_{ \theta}f(\theta_t)) - \frac{1}{n} X^Ty\\
	&= \nabla_{\theta}f(\theta_{t})- \gamma H \nabla_{\theta}f(\theta_{t})
    \end{aligned}
\end{equation}

Let us derivate $g(\gamma)=f(\alpha(\gamma))$ with respect to $\gamma$. By a composition:

\begin{equation}
    \label{eq:calcgrad}
    \begin{aligned}
	g'( \gamma) &= \langle\nabla_{\theta}f(\alpha(\gamma)),  \alpha'( \gamma) \rangle\\
	&= - \langle \nabla_{\theta}f(\theta_{t})- \gamma H \nabla_{\theta}f(\theta_t),\nabla_{\theta}f(\theta_{t}) \rangle\\
	&= - || \nabla_{ \theta}f(\theta_t) ||^2+ \gamma\langle H \nabla_{ \theta}f(\theta_t), \nabla_{ \theta}f(\theta_t)\rangle
    \end{aligned}
\end{equation}

In order to cancel the derivative, we must have that

\begin{equation}
    \gamma^* = \frac{|| \nabla_{ \theta_{t}}f ||^2}{\langle H \nabla_{\theta}f(\theta_t), \nabla_{\theta}f(\theta_t)\rangle} 
\end{equation}

We note that this is correct if $ \nabla_{ \theta}f(\theta_t)\neq 0$. If $ \nabla_{\theta}f(\theta_t)=0$, this means that $ \theta_{t}= \eta^*$, as $f$ is convex.
\\

This computation may then be done at each iteration.
\\

An important remark is that if we note $ \theta_{t+1}^* = \theta_{t}- \gamma^* \nabla_{ \theta}f(\theta_t) = \alpha(\gamma^*)$, then equations  \ref{eq:1} and \ref{eq:calcgrad} shows that

\begin{equation}
    \langle\nabla_{\theta}f(\theta_{t+1}^*), \nabla_{ \theta}f(\theta_t) \rangle=0
\end{equation}

Two optimal directions of the gradient updates are \textbf{{orthogonal}}. Importantly, this is true in the general case, not only for least-squares.

\subsubsection{\large\color{Periwinkle}Backtracking line search}

In many practical situations, it is not possible to compute explicitely the optimal step $\gamma^*$. Or it could be possible, but too expensive computationnally.

In such situations, it is possible to compute an approximation of $ \gamma^*$, for instance using \textbf{{backtracking line search.}}  This method attempts to find a good $\gamma$ by trying several decreasing values until a sufficient decrease in $f$ after the gradient update is obtained.

\url{https://en.wikipedia.org/wiki/Backtracking_line_search} 

\end{document}
